% =====================================================================
% Panel captions for the validation figures of the six catalogue papers.
%
% Every number below is read off the recorded results in
% catalogue/validation/results/*.json and the figure scripts in
% catalogue/validation/figures/.  Sub-charts are lettered A-D from left
% to right; the panel scripts tag them lowercase a-d in the rendered
% image.  One sub-chart in every panel is three-dimensional and is
% identified as such in its block.
%
% Only plain LaTeX is used -- no paper-specific \newcommand is assumed,
% so each caption is portable into any of the six preambles.  The
% \includegraphics paths are written relative to each paper's own
% directory; if this file is \input from elsewhere, prefix the paths.
% =====================================================================


% =====================================================================
% PAPER 1 --- instrument-process-ladder
% =====================================================================

\begin{figure*}[!htbp]
    \centering
    \includegraphics[width=\textwidth]{figures/panel1_composition.png}
    \caption{\textbf{The Ladder Composes Multiplicatively, and Saturation Requires a Divergent Sum.}
    (\textbf{A})~Cumulative resolving power across the rungs of a five-rung ladder with powers $[0.6, 0.35, 0.5, 0.25, 0.15]$. The composition reaches $0.9171$, and the numeric evaluation agrees with the closed form $1 - \prod_j (1 - \pi_j)$ to eleven significant figures. Composition is not an average over rungs and not a maximum over them; it is the complement of a product, which is why each additional rung buys a diminishing but never-zero increment.
    (\textbf{B})~Analytic and numeric sensitivity $\partial \pi(L) / \partial \pi_j$ against the rung power $\pi_j$, for both series superimposed. Ranking the rungs by sensitivity returns $[k_1, k_3, k_2, k_4, k_5]$, which is exactly the ordering of the powers themselves: control over the composite sits at the strongest rung and nowhere else. The two series coincide because the analytic derivative was checked against a finite-difference evaluation rather than assumed.
    (\textbf{C})~Residual $r_n$ against the number of rungs $n$ on logarithmic axes, for a divergent series of powers and a convergent one. Over $4000$ terms the divergent ladder closes to a gap of $2.50 \times 10^{-4}$ while the convergent one stalls at $0.5001$ and stays there. Saturation is a property of $\sum_i \pi_i$ diverging, not of composing many rungs: a thousand convergent rungs do not reach what a divergent tail reaches.
    (\textbf{D})~Three-dimensional surface of composite resolving power over the first two rung powers $\pi_1$ and $\pi_2$, with the remaining three held fixed. The surface is monotone and strictly concave in each argument and never leaves the unit interval---which is the property the additive and max-based rules of \textbf{Figure}~\ref{fig:ladder_cost} both lose.}
    \label{fig:ladder_composition}
\end{figure*}


\begin{figure*}[!htbp]
    \centering
    \includegraphics[width=\textwidth]{figures/panel2_contact.png}
    \caption{\textbf{Readout Tracks Contact and Nothing Else: Transit Is Free and the Product Rule Is Exact.}
    (\textbf{A})~Readout deviation against contact density. Deviation is exactly $0$ at zero density and rises monotonically to $0.007558$ at density $8$. An instrument that responded to something other than contact would show a nonzero intercept; this one does not, which is the empirical content of the claim that only contacts are read.
    (\textbf{B})~Three-dimensional surface of deviation over contact density and rung power, normalised to the value at the shortest path. Cost in contacts stays at $5$ and resolution stays at $0.917125$ as the path length is driven from $1$ to $10\,000$---a factor of $10^4$ in transit for a change of zero in both readouts. Transit between contacts is not a measurement and does not enter the ladder.
    (\textbf{C})~Mean absolute error of three composition rules against the measured composite over $n = 400$ trials. The multiplicative rule gives $0.0$ error to machine precision; the additive rule gives $0.0728$ and the max-based rule $0.1890$. The two alternatives are not slightly worse approximations of the same law---they are different laws, and only one of them is the one the ladder obeys.
    (\textbf{D})~Static verdicts against executed outcomes over $300$ cases, split into accepted and rejected. There were $0$ disagreements, and $132$ of the $300$ cases were rejected statically without being run. A compiler that rejects a ladder before execution is making a claim that execution can refute; over these cases it never did.}
    \label{fig:ladder_contact}
\end{figure*}


\begin{figure*}[!htbp]
    \centering
    \includegraphics[width=\textwidth]{figures/panel3_inertness.png}
    \caption{\textbf{Inertness: Names Do Not Move the Readout, Order Does, and Sensitivity Ranks with Power.}
    (\textbf{A})~Resolution and ambiguity under relabelling and under reordering. Over $200$ relabellings the readout moved $0$ times, while the control---a change of $0.01$ in a single rung power---separates immediately. An invariance whose control never fires is a claim that the quantity is constant, which is a weaker and different statement.
    (\textbf{B})~Three-dimensional surface of composite resolution over the first two rung powers, with the two substrates and the control marked. Both substrates read $0.917125$ to twelve digits despite differing in geometry ($3.0$ against $0.2$), field scale ($0.4$ against $7.5$) and path length ($120$ against $8$); the control, a different contact sequence, reads $0.409510$. No pair of rung powers on this surface reaches the control from the substrates' position: the difference is in the sequence, not in the parameters.
    (\textbf{C})~Rank by rung power against rank by measured sensitivity. The two orderings coincide exactly, $[k_1, k_3, k_2, k_4, k_5]$ in both. The Shapeshifter prototype compiler reproduces both the composite $0.91712$ and this ranking independently of the experiment, from the source ladder rather than from the recorded numbers.
    (\textbf{D})~Trials in which inertness held, over the $200$ relabellings. Every trial holds. The measured readout is a function of the contact sequence alone, which is what licenses treating the substrate's identity as inert.}
    \label{fig:ladder_inertness}
\end{figure*}


\begin{figure*}[!htbp]
    \centering
    \includegraphics[width=\textwidth]{figures/panel4_cost.png}
    \caption{\textbf{Cost in Contacts, Not in Transit: the Product Rule Is the Only One That Stays in the Unit Interval.}
    (\textbf{A})~Three-dimensional measured sweep of readout deviation over contact density and rung power. The surface rises in density from an exact $0$ at the origin and is flat in neither argument; every point is a measured sweep value rather than a tiled copy of a single cell.
    (\textbf{B})~Resolving power against the number of rungs composed, under the product rule as measured, the additive rule, and the max-based rule. The additive rule leaves the unit interval, which is not a matter of accuracy but of type: a resolving power above $1$ names nothing. The max-based rule stays inside but is flat, so it cannot express the gain from adding a rung at all.
    (\textbf{C})~Mean absolute error of each rule over $n = 400$ trials, plotted at the same scale as (\textbf{B}). The product rule's error is $0.0$; the other two are $0.0728$ and $0.1890$. This is the same measurement as \textbf{Figure}~\ref{fig:ladder_contact}(\textbf{C}) shown against the trajectories that produce it.
    (\textbf{D})~Share of the $300$ cases settled statically, as accepted and rejected bars. Static analysis settles $132$ of $300$ before execution with $0$ disagreements against the executed result. The ladder shortfall is reported as a compile error rather than a runtime one: a three-contact ladder against a requirement of $0.9$ compiles to \emph{``ladder shortfall: $3$ contact(s) give resolution $0.58000$, requirement is $\ge 0.9$ $\rightarrow$ UNREACHABLE''}, and the same bound run forward gives a minimum of $3$ contacts reaching $0.936$.}
    \label{fig:ladder_cost}
\end{figure*}


% =====================================================================
% PAPER 2 --- observation-groups
% =====================================================================

\begin{figure*}[!htbp]
    \centering
    \includegraphics[width=\textwidth]{figures/panel1_floor.png}
    \caption{\textbf{Grouping Is Never Free: Every Collapse Buys One Degree of Freedom and Discards a Bounded Amount.}
    (\textbf{A})~Degrees of freedom against the number of groups in the design, along a strict chain of partitions of six observations. The count runs $[0, 1, 2, 3, 4, 5]$, one degree per collapse, and is identical for two unrelated observation vectors---so the quantity is a property of the design and not of the data. At the finest partition the pooled dispersion is undefined; at the coarsest it is $1.1951$.
    (\textbf{B})~Lower bound discarded against the number of groups collapsed. The finest partition discards $0$; the coarsest discards at least $2.6355$; the sequence is monotone in between. Collapsing observations into a group is not a free simplification---it forfeits a quantity that can be bounded before the data are seen.
    (\textbf{C})~Separation against the weight of the medium. The floor is $\beta_0 = 0.5271$ and the minimum separation over the design is $5.2711$, exactly ten times it. The floor is the smallest medium weight and nothing else: deleting the medium's edges leaves a minimum cut of exactly $0.0$, so the medium does not lower the floor when removed---it destroys it.
    (\textbf{D})~Three-dimensional surface of the survivor count over the grouping lattice against the decision threshold. The lattice on six observations has $203$ elements, of which $122$ survive; the survivor set is join-closed and has a unique greatest element, which is fixed by any swap that preserves the observation vector.}
    \label{fig:groups_floor}
\end{figure*}


\begin{figure*}[!htbp]
    \centering
    \includegraphics[width=\textwidth]{figures/panel2_verdict.png}
    \caption{\textbf{The Same Data Yields Opposite Verdicts Under Different Groupings, and Pooling Shifts the Statistic by a Measurable Amount.}
    (\textbf{A})~Test statistic under two groupings of the identical observations. Grouping $P_1$ gives $6.5727$ and grouping $P_2$ gives $6.0000$; the threshold $6.2863$ falls strictly between them, so one design rejects and the other does not on numerically identical data. The verdict is a joint property of the data and the design, and reporting it without the design is reporting half of it.
    (\textbf{B})~Statistic before and after pooling, $3.5782$ falling to $3.2385$, a ratio of $1.1049$. Pooling is not a neutral aggregation: it moves the statistic by a factor that can be computed in advance from the design alone.
    (\textbf{C})~Three-dimensional measured surface of the cost of collapsing groups, over the number of groups collapsed and the family of designs. The surface is monotone in the collapse count at every design, which is the claim that no grouping recovers what a coarser one discarded.
    (\textbf{D})~Decision procedure against interval index over $26$ intervals, of which $6$ are unstable. Comparing the endpoint rule against exhaustive evaluation of every interior grouping gives $0$ disagreements: the endpoints decide the whole interval, so the exhaustive computation over the interior is unnecessary rather than merely expensive.}
    \label{fig:groups_verdict}
\end{figure*}


\begin{figure*}[!htbp]
    \centering
    \includegraphics[width=\textwidth]{figures/panel3_lattice.png}
    \caption{\textbf{The Survivor Set Is an Up-Set in the Grouping Lattice and the Floor Is Carried Entirely by the Medium.}
    (\textbf{A})~All groupings and surviving groupings against the number of groups in the partition. Of the $203$ partitions of six observations---the sixth Bell number---$122$ survive, and the survivor set is closed under join with a unique greatest element. Survivorship is therefore not a scattered property of individual designs: knowing the coarsest survivor determines the set.
    (\textbf{B})~Three-dimensional surface of separation over the medium-weight multiplier and the observation index. Separation rises with the medium at every observation and is flat in the observation index, which is the signature of a floor supplied by one vertex rather than accumulated across many.
    (\textbf{C})~Minimum cut isolating each observation. The floor $\beta_0 = 0.5271$ is attained, and the control is decisive: with the medium's edges deleted the minimum cut is exactly $0.0$ for every observation. The medium's adjacency to every other vertex is definitional in this construction, not an artefact to be pruned---a graph without it is a different theory rather than a cleaner one.
    (\textbf{D})~Pooled within-group dispersion against degrees of freedom $n - |P|$. At zero degrees of freedom the dispersion is undefined and the design computes nothing; the value rises to $1.1951$ at the coarsest partition. A design that has consumed all its degrees of freedom in grouping has no residual left to test with.}
    \label{fig:groups_lattice}
\end{figure*}


\begin{figure*}[!htbp]
    \centering
    \includegraphics[width=\textwidth]{figures/panel4_decline.png}
    \caption{\textbf{A Declined Verdict Is Informative, and a Singleton Survivor Licenses Nothing.}
    (\textbf{A})~Test statistic against cluster separation $\Delta$ for the members of one interval. The interval contains five groupings and the statistic varies across them at fixed $\Delta$, which is why the verdict must be reported against the design.
    (\textbf{B})~Members agreeing with the verdict, against cluster separation, for two intervals of sizes $5$ and $4$. Both intervals return the same verdict while disagreeing on which members support it---so agreement on a conclusion is not agreement on the reason for it, and a procedure that reports only the verdict conceals the difference.
    (\textbf{C})~Three-dimensional surface of the statistic over the interval members and $\Delta$, cut at the decision threshold. The cut plane shows the region where the verdict is stable under the design's remaining freedom; outside it the verdict is a property of the chosen grouping.
    (\textbf{D})~Interval size against surviving groupings. At one of the five groupings the verdict rests on a single survivor, with statistic $31.08$ against a next-best of $25.62$. A singleton survivor is a design in which the conclusion cannot be checked against any alternative grouping, and declining to conclude there is a stronger position than concluding from one.}
    \label{fig:groups_decline}
\end{figure*}


% =====================================================================
% PAPER 3 --- coordinate-provenance
% =====================================================================

\begin{figure*}[!htbp]
    \centering
    \includegraphics[width=\textwidth]{figures/panel1_collapse.png}
    \caption{\textbf{Collapse Is Forced by Cardinality Alone, and No Downstream Stage Recovers It.}
    (\textbf{A})~Share of coordinate maps that are auditable, against the size of the coordinate space $|\mathrm{Crd}|$. Over all $512$ maps from a three-element context space into a two-element coordinate space, exactly $0$ are auditable. The control widens the codomain to $9$ elements and an auditable witness appears immediately, so the zero is a consequence of the cardinality and not of the enumeration failing to recognise auditability.
    (\textbf{B})~Three-dimensional surface of auditability over $(|\mathrm{Crd}|, |\mathrm{Ctx}|)$. There is a zero wall wherever $|\mathrm{Crd}| < |\mathrm{Ctx}|$---a hard boundary, not a decline. The result is a counting argument, so the small case shown here is the general case.
    (\textbf{C})~Auditability restored against the $256$ second-stage maps attempted. None restores it. The control confirms the machinery is not simply reporting failure everywhere: composing an auditable first stage with an injective second stage stays auditable. Loss at stage one is not a deficiency a later stage can repair.
    (\textbf{D})~Contexts recovered against the $81$ candidate recovery functions, enumerated exhaustively after checking that the two pipeline outputs are the \emph{same element}. Zero recover the context. This holds for every function, learned or otherwise, and is not an empirical claim about a particular estimator; the control returns $9$ recoveries as soon as the inputs are made distinct.}
    \label{fig:prov_collapse}
\end{figure*}


\begin{figure*}[!htbp]
    \centering
    \includegraphics[width=\textwidth]{figures/panel2_record.png}
    \caption{\textbf{The Record Is Cheap, Its Advantage Does Not Erode with Scale, and a Fixed Record Runs Out at a Definite Release.}
    (\textbf{A})~Record size in bits against the number of measurements $N$, for the per-run scheme and the per-measurement scheme. Both are linear in $N$; only the constant differs. Over six stages with record sizes $64, 128, 32, 256, 16, 512$, per-measurement cost runs from $3900$ bits at $N = 10^2$ to $3.9 \times 10^7$ bits at $N = 10^6$.
    (\textbf{B})~Ratio of per-measurement to per-run cost across five decades of $N$. The ratio is $2.4529$ at $N = 10^2$, $2.5082$ at $10^3$, $2.5139$ at $10^4$, $2.5145$ at $10^5$ and $2.5145$ at $10^6$---flat to within $3\%$ over four decades. The control shows the statistic is capable of moving: with one very small stage the ratio reaches $6884.1$, three orders of magnitude away, so its flatness in $N$ is a measurement rather than a saturation of the instrument.
    (\textbf{C})~Three-dimensional surface of $\log_2 |\mathrm{Prv}_i|$ over the number of stages $s$ and the record size. The saving is set by the stage count and the record size and is flat in $N$ throughout, which is the whole of the reply to the cost objection.
    (\textbf{D})~Distinctions accumulating against release index, for a record of capacity $4$. A constant record is insufficient immediately; the capacity-$4$ record is exhausted at release $3$, where $6$ contexts induce $5$ behavioural classes. The floor counts distinctions, not contexts: the control adds a context that is behaviourally a duplicate and the class count stays at $2$, consuming no capacity.}
    \label{fig:prov_record}
\end{figure*}


\begin{figure*}[!htbp]
    \centering
    \includegraphics[width=\textwidth]{figures/panel3_minimality.png}
    \caption{\textbf{The Minimal Record Attains the Image-Size Bound, but the Factorisation Printed in Its Proof Is False.}
    (\textbf{A})~Sufficient record maps against the number of contexts collapsed $k$. Sufficiency thins as contexts split, and with $3$ contexts falling into $2$ behavioural classes the minimal record $\mathrm{rec}_{\min}$ is sufficient while $0$ sufficient maps have a smaller image.
    (\textbf{B})~Three-dimensional surface of the count of sufficient maps with a smaller image, over $(k, |\mathrm{Ctx}|)$. The floor is $0$ in every cell of the grid---the bound is not attained in one lucky configuration but everywhere the family was swept.
    (\textbf{C})~Record maps against cell $(|\mathrm{Ctx}|, k)$, separating the two claims the theorem makes. The image-size bound holds; the factorisation does not. Six sufficient maps fail to factor through $\mathrm{rec}_{\min}$. The counterexample is immediate once stated: the map sending $c_1, c_2, c_3$ to $0, 1, 2$ is injective and therefore sufficient, yet it splits the class $\{c_1, c_2\}$ that $\mathrm{rec}_{\min}$ identifies, and no map out of $\mathrm{rec}_{\min}$'s image can recover that split.
    (\textbf{D})~Sufficient record maps against image size. No sufficient map has an image smaller than the behavioural quotient, and the control confirms smaller maps exist and are simply not sufficient---so the minimum is a real one rather than an absence of candidates. Sufficiency is closed upwards under refinement while factorisation through a quotient requires coarseness; the two pull in opposite directions, and the universal property should be struck or restricted to sufficient maps of minimal image.}
    \label{fig:prov_minimality}
\end{figure*}


\begin{figure*}[!htbp]
    \centering
    \includegraphics[width=\textwidth]{figures/panel4_declining.png}
    \caption{\textbf{Declining Is Never Unlicensed, and a Total Comparator Buys Totality by Being Wrong Exactly Where the Partial One Declines.}
    (\textbf{A})~Admitted and declined comparisons against unsound admissions, per cell $(|\mathrm{Ctx}|, k)$. Of $45$ admitted pairs, $0$ are unsound; four context pairs are declined. The comparator is not passing its test by admitting everything.
    (\textbf{B})~Context pairs against cell, showing for each declined pair the answer a total comparator gives it. The total comparator answers $4$ unlicensed comparisons---one for each pair the partial comparator declines---while the partial comparator answers $0$ of them. The counts coincide exactly.
    (\textbf{C})~Three-dimensional surface of unlicensed answers over the family of comparators, spanned by $k$ and $|\mathrm{Ctx}|$. Declining is not a smaller answer than answering; it is the removal of precisely the answers that are not licensed, and the surface shows the removal is exact rather than conservative across the whole family.
    (\textbf{D})~Classes the record must separate against the number of contexts in the release. The floor counts distinctions rather than contexts: adding a behaviourally duplicate context leaves the class count unchanged and consumes no capacity, while a release of $6$ contexts inducing $5$ classes exhausts a capacity-$4$ record.}
    \label{fig:prov_declining}
\end{figure*}


% =====================================================================
% PAPER 4 --- runtime-graph
% =====================================================================

\begin{figure*}[!htbp]
    \centering
    \includegraphics[width=\textwidth]{figures/panel1_probe.png}
    \caption{\textbf{The Probe Is Correct at Diverging Cost, Separation Is Monotone in the Query, and the Key Is Invariant Under Relabelling.}
    (\textbf{A})~Max-flow evaluation and exhaustive minimisation over cuts returning the same answer at diverging cost, across $41$ queries. There were $0$ mismatches. The control is what makes this more than two implementations of the same mistake: perturbing the returned value by $0.17$ is detected, so the comparison has the resolution to see a disagreement of that size and reports none.
    (\textbf{B})~Three-dimensional surface of separation over the query lattice. Separation never falls as the query grows: over $150$ nested query pairs there were $0$ violations, with separations ranging from $0.7300$ to $0.9000$---so the ordering is not holding because everything is equal.
    (\textbf{C})~One probe at the top of a chain against the four probes it replaces. The chain separations are $0.9$, $0.9$, $0.9$ and $1.71$, all at or below the top value of $1.7100$, so a single probe at the top bounds all four. The control confirms the bound is a real constraint: a threshold of $0.9000$ is violated by the chain, as it must be.
    (\textbf{D})~Keys moved by relabelling against keys moved by reweighting. A weight-preserving swap moves $0$ keys; an asymmetric graph, where the swap is not weight-preserving, does move the key. An invariance whose control never fires is a claim that the key is constant, which is a different and weaker statement.}
    \label{fig:rtg_probe}
\end{figure*}


\begin{figure*}[!htbp]
    \centering
    \includegraphics[width=\textwidth]{figures/panel2_certificate.png}
    \caption{\textbf{An Answer Carries Its Own Proof: All Honest Certificates Accepted, All 123 Forgeries Rejected.}
    (\textbf{A})~Separation against query index over $56$ queries, with the floor $\beta = 0.2000$ drawn. The minimum separation is $0.7300$ and $0$ queries fall below the floor.
    (\textbf{B})~The forgery check turning over at its own tolerance. All $41$ honest certificates were accepted; against them $123$ forgeries were constructed in three families---an inflated value, a deflated value, and a certificate placing the medium inside $S$---and $0$ were accepted. The three families matter separately: a checker that verifies only the arithmetic passes the third, and one that verifies only membership passes the first two.
    (\textbf{C})~Three-dimensional surface of the guarded region over the perturbation grid. Where the guard applies it holds: with slack $0.1600$ and $m + M = 7$, $0$ of $23$ guarded perturbations moved the minimiser while $184$ unguarded ones did. The slack condition is doing the work rather than being idle.
    (\textbf{D})~The guard's width against the residual bound. Over $21$ cases the residual bound never exceeded the true slack, so the cheap bound is conservative and acting on it is safe even where it is loose. A bound that were sometimes optimistic would be unusable however tight it was on average.}
    \label{fig:rtg_certificate}
\end{figure*}


\begin{figure*}[!htbp]
    \centering
    \includegraphics[width=\textwidth]{figures/panel3_medium.png}
    \caption{\textbf{The Floor Is the Medium, Not the Contacts: Deleting It Makes Separation Free Rather Than Cheap.}
    (\textbf{A})~Separation with the medium present and with its edges deleted, for two kinds of vertex. An isolated vertex costs $0.3500$ with the medium and exactly $0.0$ without it, while a vertex with its own contacts still costs $0.5000$ without it. The two cases are easy to conflate and the measurement separates them: the medium is what makes separation cost anything for observations that have no contacts of their own.
    (\textbf{B})~Three-dimensional surface of the floor over the parameter grid, with and without the medium's edges. The lower sheet is identically zero. Deleting the medium does not lower the floor---it removes it. The medium's adjacency is definitional in this construction, and a graph that treats it as a removable artefact has a different theory rather than a cleaner one.
    (\textbf{C})~Share of keys that are shared, against the number of items $n$. At the largest size measured every key is shared, which is the boundary the paper states rather than conceals.
    (\textbf{D})~Resolution under poorer and richer weights. Richer weights raise the share of items the key resolves; the improvement is bounded, and the boundary is the subject of \textbf{Figure}~\ref{fig:rtg_stability}. The construction that defeats the key is realisable: after examining $31\,139$ pairs, two non-isomorphic graphs were found with identical keys on all $15$ queries. The rarity is the point---a sampling test would have reported the shortfall as a refutation, so the search was powered accordingly.}
    \label{fig:rtg_medium}
\end{figure*}


\begin{figure*}[!htbp]
    \centering
    \includegraphics[width=\textwidth]{figures/panel4_stability.png}
    \caption{\textbf{The Sensitivity Bound Is Attained and Never Breached, and the Cheap Residual Bound Is Conservative Everywhere.}
    (\textbf{A})~Breaches against perturbation size $\epsilon$, over $400$ perturbations at $M = 5$. There were $0$ breaches, and the worst observed move fell short of the bound by less than $10^{-5}$---so the bound is attained rather than merely respected, which distinguishes a tight bound from a vacuous one.
    (\textbf{B})~Looseness of the bound, as the ratio of the bound to the observed move, distributed across cases. The bound is loose by a roughly constant factor rather than by a factor that grows with the perturbation, so its conservatism is predictable.
    (\textbf{C})~Three-dimensional surface of the guard over the parameter grid. The guard never overstates anywhere on the grid: over $21$ cases the residual bound never exceeded the true slack.
    (\textbf{D})~Measured slack against predicted slack, one point per query. Every query sits above the diagonal, which is the direction that matters for use---a bound that is sometimes optimistic cannot be acted on, whereas one that is always conservative can, at the cost of occasionally re-running a computation that did not need it.}
    \label{fig:rtg_stability}
\end{figure*}


% =====================================================================
% PAPER 5 --- sink-detection
% =====================================================================

\begin{figure*}[!htbp]
    \centering
    \includegraphics[width=\textwidth]{figures/panel1_collapse.png}
    \caption{\textbf{A Sink Pins Every Separation into a Narrow Band, and the Bound's Depth Half Drives Nothing.}
    (\textbf{A})~Separation against sink level $\lambda$, with the floor $\beta = 0.5$ and the two-sided band drawn. At $\lambda = 2.0$ the six items separate at $2.8$, $3.1$, $3.1$, $3.1$, $3.1$, $2.8$ while the sink itself separates at $19.0$; there were $0$ band breaches and $0$ lower-bound breaches. The ceiling is attained in all $54$ cells of the $(n, \lambda)$ grid and the floor is never reached---so the upper half of the bound is exact and the lower half is slack everywhere.
    (\textbf{B})~Three-dimensional surface of the sink's weighted spread over items $n$ and $\log_{10}\lambda$. The spread rises with $\lambda$ and falls with $n$, which is the shape that lets one statistic detect a sink at any size.
    (\textbf{C})~Relative band width against sink level. The band collapses from $0.30$ at $\lambda = 0$ to $0.24$, $0.15$, $0.06$, $0.0176$ and $0.0046$ at $\lambda = 0.25, 1, 4, 16, 64$. A stronger sink makes every item's separation more nearly identical: the structure the graph was built to expose is erased, and it is erased silently, since the contaminated values ($2.8, 2.8, 3.1, 3.1, 3.1, 3.1$) are indistinguishable from a clean impersonator's.
    (\textbf{D})~Share of the separation cost carried by the sink, against $n$. The minimum share is exactly $0.6452$ at every $n$ from $3$ to $8$, while the mean drifts only from $0.6912$ to $0.6624$. The bound is flat: the proof's closing claim that ``taking $\lambda$ large or the analysis to deeper $S$ drives the fraction to $1$'' is supported for $\lambda$ and not for depth, and the depth half should be struck. The control confirms an ordinary vertex dilutes instead, from $0.048$ down to below $0.02$.}
    \label{fig:sink_collapse}
\end{figure*}


\begin{figure*}[!htbp]
    \centering
    \includegraphics[width=\textwidth]{figures/panel2_degree.png}
    \caption{\textbf{Degree Cannot Detect a Sink---the Harmless Vertex Always Has More Contacts---but Weighted Spread Can.}
    (\textbf{A})~Degree of the harmless vertex and of the fatal one, against items $n$. The harmless vertex's degree grows with $n$ from $4$ to $12$ while the sink's stays fixed at $3$. At $n = 9$ the harmless vertex has degree $9$ and the sink degree $3$: any detector that ranks by contact count puts the harmless vertex on top at every size.
    (\textbf{B})~Factor by which deleting each vertex moves a separation. Deleting the harmless degree-$9$ vertex shifts separations by at most $0.0111$ at $\epsilon = 0.1$; deleting the degree-$3$ sink multiplies them by exactly $4.0$, at every size measured. More contacts, less damage---which is why degree is not merely a weak detector here but an inverted one.
    (\textbf{C})~Three-dimensional surface showing the order reversing between the two detectors, over items $n$. Weighted spread at $n = 9$ is $0.00725$ for the harmless vertex and $0.0833$ for the sink---an order of magnitude the other way. Only one of the two detectors puts the sink on top, and it is not the one based on counting contacts.
    (\textbf{D})~The detector's ratio, fatal to harmless, against $n$. The ratio stays above unity at every size for weighted spread and below it for degree. The separation cost is a weighted quantity, so a detector that discards the weights discards exactly the information that identifies the sink.}
    \label{fig:sink_degree}
\end{figure*}


\begin{figure*}[!htbp]
    \centering
    \includegraphics[width=\textwidth]{figures/panel3_spread_sound.png}
    \caption{\textbf{The Upper Bound Holds Exactly; the Lower Bound Fails on 330 of 2403 Cases, and the Failure Does Not Thin Out with Size.}
    (\textbf{A})~Measured separation after deletion against predicted separation, one point per case, over $2403$ cases. The upper bound is breached $0$ times with a worst overshoot of $4.4 \times 10^{-16}$---machine precision, so it is exact. The lower bound is breached $330$ times.
    (\textbf{B})~Histogram of how far below the bound each breach falls, as a share of $\mathrm{sep}(v)$. The miss is a magnitude rather than a rounding error, so the failure cannot be repaired by tightening a tolerance.
    (\textbf{C})~Three-dimensional surface of $\tau$ over items $n$. Only a case whose minimiser moved leaves the plane: all $330$ breaches lie in the moved-minimiser population, and $0$ of the $1979$ stable pairs breach. The bound is sound exactly on the pairs whose minimising set is unchanged by the deletion, which is the repair the paper adopts as its corrected statement.
    (\textbf{D})~Share of breaching pairs against $n$, at $n = 3, 4, 5$: $104$ of $627$, $112$ of $716$, $114$ of $1060$. The share falls but the count does not vanish, so this is a defect in the statement rather than a small-graph artefact. A sampling test at one size would have found it and misreported its cause.}
    \label{fig:sink_spread_sound}
\end{figure*}


\begin{figure*}[!htbp]
    \centering
    \includegraphics[width=\textwidth]{figures/panel4_threshold.png}
    \caption{\textbf{A Threshold That Can Never Fire: 4776 Vertices Tested, Zero Above the Line, Best Margin $-0.0672$.}
    (\textbf{A})~Weighted spread of each vertex against the forced threshold $1 - \beta/W$. Over $4776$ vertices, $0$ lie above the line. The best margin over the whole population is $-0.0672$, at $n = 3$; it worsens with size, to $-0.1845$ at $n = 4$ and $-0.2808$ at $n = 5$. The threshold's antecedent is not merely rarely satisfied---it is unsatisfiable in this construction.
    (\textbf{B})~Three-dimensional surface of the sink's weight across $S^*$ over items $n$. The slack never goes negative: over $9808$ triples the minimum slack is $-4.4 \times 10^{-16}$, which is zero to machine precision, so $404$ cases are exactly tight and none is strictly violated. It is this identity that blocks the threshold: the same relation that makes the slack nonnegative forces the weighted spread below the threshold.
    (\textbf{C})~Threshold and weighted spread against the sink's weight $\lambda$, along the constructions in which the sink supplies the floor. Even there the gap is wide---at $\lambda = 0.005$ the threshold is $0.834$ against a spread of $0.0050$. The mechanism is not a failure to find the right regime.
    (\textbf{D})~Histogram of the margin, weighted spread minus threshold, over the $4776$ vertices. The distribution lies entirely below zero. The claim should be restated as a \emph{ranking} rule---weighted spread orders vertices by how much damage their deletion does, as \textbf{Figure}~\ref{fig:sink_degree} establishes---rather than as an absolute threshold, which no vertex in this family can cross.}
    \label{fig:sink_threshold}
\end{figure*}


% =====================================================================
% PAPER 6 --- peptide-mass-invariance
% =====================================================================

\begin{figure*}[!htbp]
    \centering
    \includegraphics[width=\textwidth]{figures/panel1_invariance.png}
    \caption{\textbf{The Cut Key Is Invariant Under Relabelling, Nearly Injective, and Bought Entirely from the Medium.}
    (\textbf{A})~Keys moved by relabelling against keys moved by a real weight change, on subgraphs of a $244$-formula NIST library. Over $288$ relabellings across six sizes a key moved $0$ times, while the control---scaling one medium edge by $0.25$---moves a key in $6$ of $6$ graphs at every size. The invariance is genuine and the statistic is not merely constant.
    (\textbf{B})~Distinct keys against the number of items. The count tracks the item count closely, rising from $3.83$ distinct keys at $n = 4$ to $9.00$ at $n = 9$. An invariant that collapses everything to one value is invariant and useless; this one separates nearly every item from every other, which is what makes the invariance worth having.
    (\textbf{C})~Three-dimensional surface of mean separation over items and medium weight, with the medium present and with its edges deleted. The lower sheet is \emph{exactly} $0.0$ in all $30$ cells---not merely cheaper. With the medium present, separation is positive in all $12$ graphs at every weight and the minimum separation rises linearly in the medium's weight ($0.2776$, $0.5551$, $1.1103$, $2.2206$, $4.4412$ at weights $0.25$ through $4$) while staying flat in the number of items.
    (\textbf{D})~Share of weights carrying several distinct minimum cuts. Over $3120$ minimum cuts, $50$ of $3069$ weights carry genuinely different cuts, so a scalar separation cannot stand in for the separating region. The share is a couple of percent and falls with size, which is why the search had to be powered rather than sampled: at some sizes a small sample finds none and would report the shortfall as a refutation.}
    \label{fig:pep_invariance}
\end{figure*}


\begin{figure*}[!htbp]
    \centering
    \includegraphics[width=\textwidth]{figures/panel2_mass.png}
    \caption{\textbf{Key and Mass Are Independent Determinations: Key-Sharing Pairs Are Grossly Mass-Distinct.}
    (\textbf{A})~Mass gaps within key-sharing pairs. Over $6000$ pairs drawn from the library, $19$ share a cut key, and none of them is near-degenerate in mass---the \emph{smallest} gap within a key-sharing pair is $6448.2$ ppm, five orders of magnitude above the tolerance at which the mass sweep bottoms out. No refinement of the mass will bring such a pair into agreement, and no sharpening of the mass will separate them by key.
    (\textbf{B})~Three-dimensional surface of ambiguity over mass tolerance and items, for both determinations. Mass ambiguity vanishes as the tolerance sharpens; key ambiguity does not. The two surfaces are not two views of one quantity.
    (\textbf{C})~The library's mass degeneracy against the swept tolerance, over $23$ points, with the recorded values marked. Mass alone resolves \emph{this} library completely, with $0$ ambiguous compounds below $9.0$ ppm, and resolution is monotone in the tolerance. That is a fact about a $244$-compound library and not a refutation of the independence claim---the two are compatible precisely because the key-sharing pairs survive it untouched, which is what independence means here and would be concealed by reporting either number alone.
    (\textbf{D})~The two determinations side by side per size. Neither reduces to the other at any size measured.}
    \label{fig:pep_mass}
\end{figure*}


\begin{figure*}[!htbp]
    \centering
    \includegraphics[width=\textwidth]{figures/panel3_closure.png}
    \caption{\textbf{Closure Names a Region, Not Always an Item, and the Excluded Regime Is What Makes It a Singleton.}
    (\textbf{A})~Share of closed determinations naming a region rather than a single item, against size. Over $1560$ closed determinations, $51$ name a plural region. The reading on which closure always names one item is denied in principle by the paper's own strengthened theorem, whose proof exhibits a closed determination with a two-element ambiguity set; the measurement finds that case at every scale.
    (\textbf{B})~Mean and largest region size against items. The mean sits just above one at every size and the largest observed region is three. One swept size found no plural region at all---mean exactly $1.0$ over $280$ determinations---while both its neighbours exceed it, which makes that cell a sampling gap rather than a structural fact.
    (\textbf{C})~Three-dimensional surface of mean $|S^*|$ over contact scale and items. While no item--item contact exceeds the medium, the mean size of the minimising set is exactly $1.0$ to machine precision, in all $30$ cells at scale $\le 1$, at every size. This is what the non-completeness normalisation buys.
    (\textbf{D})~The same transition in section, with the excluded regime shaded. Past the medium the floor breaks and the minimiser grows with the number of items, reaching $5.125$ at contact scale $3$ (against $1.0$, $1.0$, $1.0$, $1.0$, $1.0$, then $1.5208$, $3.8125$, $5.125$ across scales $0.1$ to $3.0$). The excluded regime is doing work rather than being a convenience: the axiom is what separates a singleton minimiser from one that grows without bound.}
    \label{fig:pep_closure}
\end{figure*}


\begin{figure*}[!htbp]
    \centering
    \includegraphics[width=\textwidth]{figures/panel4_three.png}
    \caption{\textbf{The Only Claim the Suite Did Not Reproduce, and the Fault Is in the Statement Rather Than the Proof.}
    (\textbf{A})~The two readings of the three-catalyst proposition against the number of catalysts $k$. Read literally, \emph{``the erroneous catalyst''} is definite---identified, removed, majority taken over the survivors---and $k = 2$ already suffices, recovering the true region in $1.000$ of runs. Read as the proof's reason describes it, the erroneous catalyst is \emph{not} identified and the majority runs over all $k$; there $k = 2$ recovers the true region in $0.000$ of runs and $k \ge 3$ is required.
    (\textbf{B})~Gap between the two readings, per candidate count. The gap is exactly $1$ at $k = 2$ and exactly $0$ at every other $k$. That single point is the whole of the disagreement.
    (\textbf{C})~Three-dimensional surface of both readings over $(k, \text{candidates})$. The disagreement is identical at $3$, $4$ and $5$ candidates, so it is not an artefact of the candidate pool but of the wording. Each recorded point rests on $2000$ trials.
    (\textbf{D})~The swept curves against the recorded values, which they pass through exactly. The proposition should read \emph{``an erroneous catalyst''} and state explicitly that the erroneous one is unidentified; with that amendment the bound $k \ge 3$ is correct and tight, since at $k = 3$ both readings recover the true region in $1.000$ of runs. It is reported here as not reproduced rather than silently repaired.}
    \label{fig:pep_three}
\end{figure*}
